\documentclass[12pt]{article}
\usepackage{amsmath,amssymb,amsthm}
\usepackage[margin=1in]{geometry}
\usepackage{algorithm}
\usepackage{algpseudocode}

\newtheorem{theorem}{Theorem}
\newtheorem{lemma}[theorem]{Lemma}

\title{Problem 202: Laserbeam}
\author{Project Euler Solution}
\date{}

\begin{document}
\maketitle

\section{Problem Statement}
Three mirrors form an equilateral triangle with reflective inner surfaces and infinitesimal
gaps at vertices. A laser beam enters vertex~$C$, bounces off exactly $R = 12017639147$
surfaces, then exits through~$C$. How many distinct beam paths exist?

\section{Mathematical Foundation}

\begin{theorem}[Unfolding Principle]
A laser beam making $R$ reflections inside an equilateral triangle corresponds to a
straight line in the triangular tessellation of the plane, from the origin to a lattice
point $(a,b)$ with $a + b = n$, where $n = (R+3)/2$.
\end{theorem}

\begin{proof}
Each reflection is equivalent to reflecting the triangle across the mirror wall and
continuing in a straight line. The beam crosses $R$ edges in the tessellation. Accounting
for the entry and exit through vertex gaps, the endpoint satisfies
$a + b = (R+3)/2$. For $R = 12017639147$, we get $n = 6008819575$.
\end{proof}

\begin{theorem}[Exit Vertex Classification]
The beam exits through vertex~$C$ if and only if $b \equiv 2 \pmod{3}$, given that
$n \equiv 1 \pmod{3}$.
\end{theorem}

\begin{proof}
In the tessellation, a lattice point $(a,b)$ with $a+b = n$ is a copy of vertex~$C$ precisely when
$b \not\equiv 0 \pmod{3}$ and $a \not\equiv 0 \pmod{3}$. Since $a = n - b$ and
$n \equiv 1 \pmod{3}$, we have $a \equiv 1 - b \pmod{3}$. The condition
$a \not\equiv 0$ becomes $b \not\equiv 1 \pmod{3}$. Combined with $b \not\equiv 0 \pmod{3}$,
this forces $b \equiv 2 \pmod{3}$.
\end{proof}

\begin{theorem}[No Intermediate Vertex]
The beam passes through no intermediate vertex if and only if $\gcd(b, n) = 1$.
\end{theorem}

\begin{proof}
If $d = \gcd(a,b) > 1$, the point $(a/d, b/d)$ is an intermediate lattice point on the segment,
corresponding to a vertex through which the beam would exit prematurely. Since $a = n - b$,
we have $\gcd(a,b) = \gcd(n,b)$.
\end{proof}

\begin{theorem}[M\"obius Inversion Count]
The number of valid paths is
\[
\mathrm{Count} = \sum_{d \mid n} \mu(d) \cdot C(d),
\]
where $C(d) = \#\{b \in [1, n-1] : d \mid b,\; b \equiv 2 \pmod{3}\}$.
\end{theorem}

\begin{proof}
By M\"obius inversion, $[\gcd(b,n)=1] = \sum_{d \mid \gcd(b,n)} \mu(d)$. Exchanging
summation order:
\[
\mathrm{Count} = \sum_{d \mid n} \mu(d) \sum_{\substack{b=1 \\ d \mid b,\; b \equiv 2\!\!\pmod{3}}}^{n-1} 1
= \sum_{d \mid n} \mu(d) \cdot C(d).
\]
Substituting $b = dk$, the inner sum becomes $\#\{k \in [1, n/d - 1] : k \equiv 2d^{-1}\!\!\pmod{3}\}$,
which evaluates to $\lfloor(n/d - 1 - r)/3\rfloor + 1$ with $r = 2d^{-1}\!\!\bmod 3$.
\end{proof}

\begin{lemma}[Factorization]
$n = 6008819575 = 5^2 \times 11 \times 17 \times 23 \times 29 \times 41 \times 47.$
The 7 distinct primes yield $2^7 = 128$ squarefree divisors.
\end{lemma}

\begin{proof}
Verified by trial division. Since $\mu(d) = 0$ for $d$ divisible by any $p^2$, we enumerate
only squarefree divisors formed from $\{5, 11, 17, 23, 29, 41, 47\}$.
\end{proof}

\section{Algorithm}

\begin{algorithmic}[1]
\State $n \gets (R + 3)/2 = 6008819575$
\State Factor $n$: primes $= [5, 11, 17, 23, 29, 41, 47]$
\State $\mathrm{count} \gets 0$
\For{each subset $T \subseteq \mathrm{primes}$}
    \State $d \gets \prod_{p \in T} p$, \quad $\mu_d \gets (-1)^{|T|}$
    \State $m \gets n/d$, \quad $r \gets 2 \cdot d^{-1} \bmod 3$
    \State $C_d \gets \lfloor(m - 1 - r)/3\rfloor + 1$
    \State $\mathrm{count} \gets \mathrm{count} + \mu_d \cdot C_d$
\EndFor
\State \Return count
\end{algorithmic}

\section{Complexity Analysis}
\begin{itemize}
    \item \textbf{Time:} $O(\sqrt{n})$ for factorization, $O(2^k) = O(128)$ for the M\"obius sum.
    \item \textbf{Space:} $O(k) = O(7)$.
\end{itemize}

\section{Answer}

\[
\boxed{1209002624}
\]

\end{document}
